% titel: Blatt 3, A4: Ausbeute (Nitrobenzol aus Benzol)
% skript: Mortimer Kap. 4.5 (Ausbeute, Stöchiometrie)

% Teil I: Werkzeuge, kein Lösungsweg.

Die tatsächlich erhaltene Produktmenge liegt oft unter der theoretisch
möglichen. Aus vorliegender (realer) Produktmenge und Ausbeute lässt sich
zurückrechnen, wie viel Edukt einzusetzen ist.

\begin{werkzeugkasten}[Ausbeute]
\begin{equation}
  \eta = \frac{n_{\text{real}}}{n_{\text{theor}}}
  \qquad(\text{bzw.\ auch mit Massen, wenn } M \text{ konstant}).
\end{equation}
$\eta = \qty{76.3}{\percent}$ bedeutet: es fällt nur
$0.763\cdot n_{\text{theor}}$ an. Für eine Zielmenge
$n_{\text{real}}$ gilt also
\begin{equation}
  n_{\text{theor}} = \frac{n_{\text{real}}}{\eta}.
\end{equation}
Diesen theoretischen Wert über die Stöchiometrie (1:1 bei dieser Reaktion) in
die benötigte Edukt-Stoffmenge und Masse umrechnen.

\textit{Skript:} Mortimer Kap. 4 (Ausbeute), Übung 3/4.
\end{werkzeugkasten}

\begin{achtung}[Darauf achten]
\begin{itemize}
  \item Ausbeute als Bruch einsetzen ($0.763$), nicht als Prozentzahl
        ($76.3$).
  \item Erste Umrechnung: Zielmasse $\rightarrow$ reale Stoffmenge;
        die realen Werte beziehen sich auf das \emph{tatsächlich} erhaltene
        Produkt.
  \item Die Gleichung Benzol $+$ $\mathrm{HNO_3}$ ist hier im Verhältnis 1:1;
        bei anderen Aufgaben den Stöchiometriefaktor prüfen.
  \item Am Schluss Masse der Edukte $>$ Masse des erhaltenen Produkts, weil
        die Ausbeute unter 100\% liegt --- als Plausibilitätscheck.
\end{itemize}
\end{achtung}
